% ============================================================
% CUMCM 国赛 B 题优化类模板 — v7.2
% 编译:xelatex main.tex(跑两遍)
% 金标准:六段子结构 / 三段式公式 / 四重检验 / 算法对比表 / 创新量化
% 核心强化:目标函数+约束 / 多算法对比 / 灵敏度
% ============================================================
\documentclass[12pt,a4paper]{ctexart}

\usepackage{amsmath,amssymb,amsfonts}
\usepackage{graphicx}
\usepackage{booktabs}
\usepackage{geometry}
\usepackage{float}
\usepackage{caption}
\usepackage{subcaption}
\usepackage{listings}
\usepackage{hyperref}
\usepackage{enumitem}
\usepackage{xcolor}
\usepackage{algorithm}
\usepackage{algpseudocode}
\usepackage{indentfirst}   % 首行缩进
\usepackage{titlesec}      % 章节标题格式
\usepackage{fancyhdr}      % 页眉页脚
\usepackage{setspace}      % 行距

\geometry{top=2.5cm, bottom=2.5cm, left=2.5cm, right=2.5cm, headheight=14pt}

% ============ 排版规范（对齐国赛优秀论文 B159） ============
% 行距 1.5 倍，首行缩进 2 字符
\linespread{1.5}\selectfont
\setlength{\parindent}{2em}

% 章节标题居中，二级标题左对齐
\titleformat{\section}{\centering\large\bfseries\heiti}{\thesection}{0.8em}{}
\titleformat{\subsection}{\normalsize\bfseries\heiti}{\thesubsection}{0.8em}{}
\titleformat{\subsubsection}{\normalsize\bfseries\heiti}{\thesubsubsection}{0.6em}{}

% 页眉空、页脚页码居中
\pagestyle{fancy}
\fancyhf{}
\fancyfoot[C]{\small\thepage}
\renewcommand{\headrulewidth}{0pt}

\setCJKmainfont{SimSun}[BoldFont=SimHei, ItalicFont=KaiTi]
\setCJKsansfont{SimHei}
\setCJKmonofont{FangSong}

\DeclareMathOperator{\erf}{erf}
\DeclareMathOperator{\sgn}{sgn}

\hypersetup{colorlinks=true, linkcolor=black, citecolor=black, urlcolor=blue}

\lstset{
  basicstyle=\ttfamily\small,
  breaklines=true,
  frame=single,
  numbers=left,
  numberstyle=\tiny,
  keywordstyle=\color{blue},
  commentstyle=\color{green!50!black},
}

\title{\textbf{[研究对象]的[优化目标]建模与求解}}
\author{队伍编号:\textbf{XXXX}}
\date{}

\begin{document}

\maketitle

% ============================================================
% 摘要三段式
% ============================================================
\begin{abstract}
\noindent
\textbf{第一段}:行业背景 + 研究对象 + 三层框架(建模-求解-检验)。

\textbf{第二段}:问题一建立[模型],采用 SA-PSO 求解,年均功率 $54.2$ MW,较传统布局提升 $19.6\%$;问题二...;问题三...

\textbf{第三段}:3 项创新(精度提升 $12.3\%$ / 收敛速度提升 $41.2\%$ / 计算效率提升 $45\%$)+ 四重检验指标 + 跨行业迁移。

\vspace{0.5em}
\noindent\textbf{关键词}:[对象] [模型] [算法] [方法] [指标]
\end{abstract}

\newpage

% ============================================================
\section{问题重述}
\subsection{问题背景}
[背景]
\subsection{问题提出}
\begin{itemize}
  \item 问题一:[布局优化]
  \item 问题二:[参数寻优]
  \item 问题三:[方案对比]
\end{itemize}

% ============================================================
\section{问题分析}
[全局建模流程图]
\begin{figure}[H]
\centering
\includegraphics[width=0.85\textwidth]{../figures/pdf/fig1_roadmap.pdf}
\caption{图1 全局建模流程(六模块:数据-建模-求解-检验-对比-交付)}
\label{fig:roadmap}
\end{figure}

% ============================================================
\section{模型假设}
\begin{enumerate}
  \item 假设1:[内容](量化误差 $\leq X\%$)
  \item 假设2:[内容]
  \item 假设3:[内容]
  \item 假设4:[内容]
  \item 假设5:[内容]
  \item 假设6:[内容]
\end{enumerate}

% ============================================================
\section{符号说明}
\begin{table}[H]
\centering
\caption{决策变量与参数说明}
\label{tab:symbols}
\begin{tabular}{lllll}
\toprule
符号 & 含义 & 取值范围 & 单位 & 类型 \\
\midrule
$r_i$ & 第 $i$ 镜距塔距离 & $[50, 350]$ & m & 连续 \\
$\theta_i$ & 第 $i$ 镜方位角 & $[0, 2\pi]$ & rad & 连续 \\
$h$ & 吸热塔高度 & $[80, 200]$ & m & 连续 \\
$w$ & 镜面宽度 & $[4, 12]$ & m & 离散 \\
\bottomrule
\end{tabular}
\end{table}

% ============================================================
\section{模型建立与求解}

\subsection{问题一}

\subsubsection{参数配置与数据预处理}
[内容]

\subsubsection{基础经典模型}
年均功率 $P_{\text{avg}}$ 综合反映镜场效率与吸热塔匹配度,优化目标是最大化年均输出功率。
\begin{equation}
\max_{\mathbf{x}} \; P_{\text{avg}} = \frac{1}{8760} \sum_{t=1}^{8760} \eta_{\cos}(t) \cdot \eta_{\text{atm}}(t) \cdot I_{\text{DNI}}(t)
\label{eq:obj}
\end{equation}
其中 $\mathbf{x} = (r_i, \theta_i, h)$ 为决策变量。目标函数为年化功率最大化,需满足下列四类约束。

\paragraph{约束条件}
\textbf{(1) 几何约束}
\begin{align}
r_i &\geq r_{\min} = 50 \text{ m}, \quad \forall i \label{eq:cons1}\\
h_{\min} \leq h &\leq h_{\max} \label{eq:cons2}
\end{align}

式~\eqref{eq:cons1} 和式~\eqref{eq:cons2} 限定镜面几何范围。

\textbf{(2) 物理约束}
\begin{equation}
\eta_{\text{shadow}}(i,j) \leq \epsilon_{\text{tol}}, \quad \forall i \neq j
\label{eq:cons3}
\end{equation}

式~\eqref{eq:cons3} 约束阴影遮挡。

\textbf{(3) 工程约束}
\begin{equation}
\sum_{i=1}^{N} A_i \leq A_{\text{total}} = 1000 \text{ m}^2
\label{eq:cons4}
\end{equation}

式~\eqref{eq:cons4} 限制总镜面面积不超过 $1000\text{ m}^2$。

\subsubsection{传统模型缺陷剖析}
传统均匀布局存在以下缺陷:
\begin{enumerate}
  \item 缺陷1:未考虑三维阴影遮挡,效率损失 $X\%$
  \item 缺陷2:固定步长排布,土地利用率仅 $Y\%$
  \item 缺陷3:缺乏全局优化,易陷入局部最优
\end{enumerate}

\subsubsection{创新改进模型}
本文建立基于三维阴影遮挡的效率计算模型,并提出 SA-PSO 混合优化算法。相比式~\eqref{eq:obj} 的传统模型,精度提升 $12.3\%$。

\subsubsection{算法设计与求解}

\paragraph{三算法对比表(强制)}
\begin{table}[H]
\centering
\caption{三算法性能对比}
\label{tab:algo_compare}
\begin{tabular}{lccccc}
\toprule
算法 & 收敛迭代 & 最优值 & 稳定性 (CV) & 时间 (s) & 评级 \\
\midrule
SA-PSO & 127 & 54.2 MW & 2.3\% & 186 & ★★★ \\
GA & 215 & 52.8 MW & 5.1\% & 245 & ★★☆ \\
DE & 168 & 53.5 MW & 3.8\% & 203 & ★★☆ \\
\bottomrule
\end{tabular}
\end{table}

选用 SA-PSO 的理由:收敛最快(127 代),稳定性最高(CV=2.3\%),且最优值最大。

\paragraph{SA-PSO 伪代码}
\begin{algorithm}[H]
\caption{SA-PSO 混合优化算法}
\begin{algorithmic}[1]
\Require 目标函数 $f$, 约束 $g$, 种群大小 $N$, 最大迭代 $T_{\max}$
\Ensure 最优解 $\mathbf{x}^*$, 最优值 $f(\mathbf{x}^*)$
\State 初始化种群 $\{\mathbf{x}_1, \ldots, \mathbf{x}_N\}$
\State 计算适应度,初始化 $\mathbf{p}_i$ 和 $\mathbf{g}$
\For{$t = 1$ to $T_{\max}$}
  \State 温度退火 $T_t = T_0 \cdot \alpha^t$
  \For{每个粒子 $i$}
    \State 更新速度 $\mathbf{v}_i \leftarrow w\mathbf{v}_i + c_1 r_1 (\mathbf{p}_i - \mathbf{x}_i) + c_2 r_2 (\mathbf{g} - \mathbf{x}_i)$
    \State SA 扰动 $\mathbf{x}_i \leftarrow \mathbf{x}_i + \mathbf{v}_i + \eta \cdot T_t$
    \State 约束处理:违反约束则投影回可行域
    \State 更新 $\mathbf{p}_i$, $\mathbf{g}$
  \EndFor
  \State 若收敛则 break
\EndFor
\State \Return $\mathbf{g}, f(\mathbf{g})$
\end{algorithmic}
\end{algorithm}

\subsubsection{数值结果与可视化}
\begin{figure}[H]
\centering
\includegraphics[width=0.85\textwidth]{../figures/pdf/fig2_convergence.pdf}
\caption{图2 SA-PSO 收敛速度较 PSO 提升 $41.2\%$}
\label{fig:convergence}
\end{figure}

\begin{table}[H]
\centering
\caption{问题一优化结果对比}
\label{tab:result1}
\begin{tabular}{lccc}
\toprule
指标 & 传统布局 & 优化布局 & 提升幅度 \\
\midrule
年均功率 (MW) & 45.3 & 54.2 & +19.6\% \\
峰值功率 (MW) & 62.1 & 68.5 & +10.3\% \\
镜面利用率 (\%) & 78.2 & 89.7 & +11.5\% \\
计算时间 (s) & - & 186 & - \\
\bottomrule
\end{tabular}
\end{table}

图~\ref{fig:convergence} 展示了 SA-PSO 与 GA 的收敛曲线对比。SA-PSO 在第 127 代收敛,较 GA(215 代)提升 $41\%$。最优功率 $54.2$ MW 较 GA($52.8$ MW)提高 $2.6\%$。

\subsection{问题二}
% 同六段子结构

\subsection{问题三}
% 同六段子结构

% ============================================================
\section{模型检验}

\subsection{拟合精度检验}
$R^2 = 0.987$, $\text{MAE} = 0.23$, $\text{RMSE} = 0.31$, $\text{MAPE} = 2.3\%$。

\subsection{灵敏度分析}
\begin{table}[H]
\centering
\caption{参数灵敏度分析}
\label{tab:sensitivity}
\begin{tabular}{lccccc}
\toprule
参数 & 基准值 & ±10\% 影响 & ±20\% 影响 & 弹性系数 $\varepsilon$ & 分级 \\
\midrule
镜宽 $w$ & 6 m & ±1.2\% & ±2.8\% & 0.14 & 低 \\
塔高 $h$ & 120 m & ±3.5\% & ±7.8\% & 0.39 & 中 \\
镜数 $N$ & 200 & ±5.2\% & ±11.3\% & 0.57 & 高 \\
\bottomrule
\end{tabular}
\end{table}

\begin{figure}[H]
\centering
\includegraphics[width=0.85\textwidth]{../figures/pdf/figN_tornado.pdf}
\caption{图N 灵敏度龙卷风图($N$ 为高灵敏度参数)}
\label{fig:tornado}
\end{figure}

\subsection{蒙特卡洛验证}
$N_{\text{sim}} = 200$, $\mu = 53.8$ MW, $\sigma = 1.2$ MW, $\text{CV} = 2.2\%$, $95\%\text{CI} = [51.4, 56.2]$。

\subsection{假设误差量化}
\begin{table}[H]
\centering
\caption{假设误差量化}
\label{tab:assumption_error}
\begin{tabular}{lcc}
\toprule
假设 & 偏差量化 & 对结果影响 \\
\midrule
假设1:镜面理想平面 & $\delta = 0.5\%$ & 功率损失 $1.2\%$ \\
假设2:无大气衰减 & $\delta = 2.1\%$ & 功率损失 $3.5\%$ \\
假设3:稳态光照 & $\delta = 1.8\%$ & 功率波动 $\pm 2.0\%$ \\
\bottomrule
\end{tabular}
\end{table}

% ============================================================
\section{模型优缺点与改进}
\textbf{创新点}:
\begin{enumerate}
  \item 建立基于三维阴影遮挡的效率计算模型,精度提升 $12.3\%$
  \item 提出 SA-PSO 混合优化算法,收敛速度提升 $41.2\%$,稳定性提高 $35\%$
  \item 构建多目标优化框架,帕累托前沿求解效率提升 $28\%$
\end{enumerate}

% ============================================================
\section{参考文献}
\begin{thebibliography}{99}
\bibitem{ref1} 作者. 标题[J]. 期刊, 年份, 卷(期): 页码.
\bibitem{ref2} Author A. Title[J]. Journal, Year, Vol: Pages.
% 近5年 ≥40%, 外文 ≥30%, 总数 ≥10
\end{thebibliography}

% ============================================================
\appendix
\section{附录1 数据说明}
[result 样表 5 行 8 列]

\section{附录2 代码}
\begin{lstlisting}[language=Python, caption=problem1_optimize.py]
from sa_pso import SA_PSO

def objective(x):
    """目标函数:最大化年均功率"""
    pass

def constraints(x):
    """约束条件"""
    pass

solver = SA_PSO(obj=objective, dim=10, bounds=[(0, 100)] * 10, constraints=constraints)
result = solver.solve()  # {'x_opt', 'f_opt', 'history', 'iterations'}
\end{lstlisting}

\section{附录3 公式推导}
[目标函数完整推导]

\end{document}
